\begin{table}[htbp]
    \centering
    \begin{tabular}{ |p{2.5cm}||p{2.5cm}|p{2.5cm}|p{2.5cm}|p{2.5cm}| }
        \hline
        \hline
        Filename & Species & Control Points DPI & Control Points (mm) & Control Point Distance (Pixels) \\
        \hline
        core-scan-ac-sa-scalebar.tif & \textit{Acer saccharum} & 21145 & 101.3 & 84332.1 \\ % 84332.13 pixels or 21145 DPI
        core-scan-be-ni-scalebar.tif & \textit{Betula nigra} &  21108 & 100.8 & 83769.9 \\ % 83769.86 pixels or 21108 DPI
        core-scan-ps-me-scalebar.tif & \textit{Pseudotsuga menziesii} &  21100 & 101.13 & 84011.1 \\ % 84011.05 or 21100 DPI
        core-scan-qu-ru-scalebar.tif & \textit{Quercus rubra} &  21169 & 40.3 & 84200.1 \\ % 84200.08 or 21168.8
        core-scan-ts-me-scalebar.tif & \textit{Tsuga mertensiana}  & 21097 & 101.14 & 83998.1 \\ % 83998.05 or 21097.1
        stage-micrometer-21080.tif & N/A & N/A & N/A & N/A \\
        \hline
    \end{tabular}
    \caption{Two control points were added to these scans at approximately the same distance with an ultra-fine tipped felt pen. The distance between the two points can be measured physically with digital calipers and digitally with an image annotation tool. The DPI measurements are similar to performing the same measurements on a stage micrometer at the same lens zoom. Using the stage micrometer to measure DPI is preferred as finding the centerpoint of The scans are stored in a Zenodo repository (https://doi.org/10.5281/zenodo.22215691).}
    \label{tab:scans_control_points}
\end{table}